% 本章共同的输出评价；原连续参数风险模板已按本轮模型讲解取舍移除。
各类回归输出在拟合、选择、校准与最终测试中的分工见表\ref{tab:regression-evaluation-protocol}。

\begin{table*}[htbp]
  \centering
  \small
  \begin{tabularx}{\linewidth}{@{}>{\RaggedRight\arraybackslash}p{22mm}>{\RaggedRight\arraybackslash}X>{\RaggedRight\arraybackslash}X>{\RaggedRight\arraybackslash}X@{}}
    \toprule
    输出 & 拟合与验证选择 & 独立校准 & 最终测试 \\
    \midrule
    条件均值 & 平方损失或均值模型似然；验证MSE & 点预测不另设 & MSE、残差与分组误差 \\
    条件分位数 & 对应$\tau$的损失；检查尾部记录与交叉 & 原始分位点不另设 & 分位数损失与经验超越率 \\
    模型内响应预测区间 & 拟合响应或后验模型；选择超参数 & 按模型构造；若再校准则单独留出 & 响应覆盖率、宽度与分组表现 \\
    共形覆盖区间 & 拟合基础预测器；校准前选定模型 & 留出记录估计不合格分数分位数 & 覆盖率、宽度与分组表现 \\
    完整分布 & 似然或适当评分；验证分布诊断 & 若作再校准则单独留出 & 负对数得分或CRPS、校准与锐度 \\
    \bottomrule
  \end{tabularx}
  \caption{不同回归输出的评价分工。两类区间均以未来响应为覆盖对象：模型内响应预测区间依赖其响应模型与参数估计或后验设定，共形覆盖区间用独立校准记录估计覆盖修正。最终测试留到模型选择与校准结束之后。}
  \label{tab:regression-evaluation-protocol}
\end{table*}

评价应与交付对象对应。在逐个比较指标之前，先要保证一件事：用来报告最终水平的数据，不能在模型训练或挑选时就被看过。因此一条记录在整个流程中要分工明确——拟合与选择超参数用训练折和验证折，共形区间另用独立校准记录，最后才用完全没参与任何选择的测试集报告结果；表\ref{tab:regression-evaluation-protocol}把这套分工按交付对象列清。在此基础上，条件均值比较测试均方误差（MSE），即残差平方的平均；分位数比较相应损失及经验超越率；未来响应区间同时报告覆盖率、宽度和分组表现。参数区间与潜在函数值区间覆盖各自的真值，不能直接拿带噪响应当条件均值真值。GP预测响应时应使用$v_\star+\sigma^2$，而不是只用函数方差$v_\star$。

完整分布可用\term{适当评分}（proper scoring rule）比较：在相关期望存在时，报告真实分布能取得最优期望评分。统一按损失越小越好计，负对数得分为$-\log p(y)$；给实际响应分配的概率或密度越低，损失越大。连续密度可以超过$1$，因此该得分不必非负。连续秩概率得分（continuous ranked probability score, CRPS）比较累积分布$F$与观测所给的阶跃函数：
\begin{equation}
 \operatorname{CRPS}(F,y)=\int_{\mathbb R}
        [F(t)-\mathbf1\{y\leq t\}]^2\,dt.
 \label{eq:regression-crps}
\end{equation}
有限一阶矩保证此积分有限，并保留响应单位；若$F$只在单点$a$取值，两个阶跃函数仅在$a$和$y$之间不同，积分就等于$|a-y|$。同均值、不同分散程度的分布可以获得不同评分，均值MSE不能替代整条分布的评价\scite{gneiting2007scoring}。

在房价任务中，先按部署时间或地区划分数据，每个训练折内完成缺失值处理、标准化、特征选择与拟合，再用验证结果选择表示和超参数。若需要共形覆盖区间，应在模型冻结后用独立校准记录确定修正，最后才用未参与选择或校准的测试集报告结果。看过测试结果后再调模型，该批数据便不再是最终测试证据。模型内区间是否需要另作校准，取决于交付要求，不由“报告区间”自动决定。
